% page 378 of the facsimile (image p-377 of the scan)
% block 160.0 x 237.7 mm ; left margin 8.0 mm
\pgc{-12.700mm}{0.000mm}{
\l{\sou{metopo}. (\sou{arkitekturo}). Interspaco quadrata, maxim-multa-}
\l{\cel{3}kaze ornita per skulturi - qua esas inter la triglifi}
\l{\cel{3}di la frizo dorika. - DEFIRS.}
\l{}
\l{metro. I. Unajo di longeso qua reprezentas la 1/10.000.000-}
\l{\cel{3}ima parto di la quarimo di la meridiano Terglobala. -}
\l{\cel{3}II. Nombro e dispozeso di la silabi ye qui konsistas sin-}
\l{\cel{3}gla verso, singla versajo, poemeto, e c. - DEFIRS.}
\l{}
\l{\sou{metriko}. Ensemblo ek la reguli qui koncernas la metro di la}
\l{\cel{3}versi. -}
\l{}
\l{\sou{metrito}. (\sou{patol}.)Inflamuro di la utero. F.}
\l{}
\l{\sou{metrologio}. La cienco di pondero e mezuro.}
\l{}
\l{\sou{metronomo}. (\sou{muziko)}. Instrumento, uzata por indikar la rapi-}
\l{\cel{3}deso-gradi diversa di la movimento muzikala, per pendulo}
\l{\cel{3}qua skandas plu o min rapide segun ke onu plualtigas o}
\l{\cel{3}minaltigas pezajo movebla sur la stangeto gradizita.- DEFIRS}
\l{}
\l{\sou{metropolo}. I. Urbo, stato, konsiderata kompare a la kolonii}
\l{\cel{3}qui naskis de lu. - II. (\sou{epoki antiqua}). Urbo chefala di}
\l{\cel{3}provinco. - III. (historio religiala). Urbo qua havas re-}
\l{\cel{3}zideyo arkiepiskopala. - DEFIRS.}
\l{}
\l{\sou{metropolito}. (\sou{eklezio Greka}). Granda oficiero di la eklezio}
\l{\cel{3}ortodoxa, di qua la rango hierarkiala esas inter la patri-}
\l{\cel{3}arko e la arkiepiskopi. - DEFIRS.}
\l{}
\l{\sou{mezo}. Parto qua distas egale de la extremaji, en intervalo}
\l{\cel{3}de loki o de epoki. - I.}
\l{}
\l{\sou{mezentero}. (\sou{anat.}) Rifalduro di la peritoneo, qua mantenas}
\l{\cel{3}la parti diversa di la intestini. - DEFIS.}
\l{}
\l{\sou{mezereono}. (\sou{bot.}) Arbusto ek la genero \textquotedbl{}dafnei\textquotedbl{}, qua kreskas}
\l{\cel{3}sur la monti arboroza, boskoza, di Francia e di la centro-}
\l{\cel{3}regioni di Europa. - L.\cel{1}\sou{daphne mezereum}.}
\l{}
\l{\sou{mezodermo}. I. (\sou{embriol.})Strato qua aparas, dum la su-developo}
\l{\cel{3}prima, inter la du folieti di ektodermo ed endodermo.-}
\l{\cel{3}II. (\sou{bot.})\cel{2}Parto di la kortico, inter la strato korkatra}
\l{\cel{3}e la envelopilo herbatra. -DEFIS.}
\l{}
\l{\sou{mezokarpo}. (\sou{bot.)}\cel{3}Strato di la perikarpo di frukto, inter}
\l{\cel{3}la epidermo extera (epikarpo) e la epidermo interna (endo-}
\l{\cel{3}karpo) - pulpoza, che ula frukti (cerizi, abrikoti, e c.)}
\l{\cel{3}- lignatra che altrai (bruo). - DEFI S.}
\l{}
\l{mezolabo. Olima instrumento di matematiki, qua konsistas y e}
\l{\cel{3}tri paralelogrami, moviva sur kuliso, uzita por trovar}
\l{\cel{3}mekanikale du meza proporcionali quin onu ne povis deter-}
\l{\cel{3}minar geometriale en la problemo \textquotedbl{}duopligo di la kubo\textquotedbl{}.}
\l{\cel{3}- DEFIRS.}
}
